% ============================================================================
% Fichier cible : ARIEL_LOGMINER_MEMOIRE_FINAL/chapters/chapitre2_etat_art.tex
% Passage à remplacer : de « \section{Mise à jour et gouvernance des modèles} »
%   (inclus) à « \section{Tableaux de bord et exploitabilité par l'analyste} »
%   (exclu).
% Agent : Agent 3 — État de l'art et positionnement, chantier 2026-11-05.
% Sources : Sculley et al. ; Breck et al. ; Amershi et al. ; NIST AI RMF 1.0.
% ============================================================================

\section{Gouvernance des modèles et MLOps}

\subsection{Le modèle comme composant d'un système évolutif}

Le déploiement d'un modèle ne se résume pas à sérialiser un estimateur. Sculley et al.
montrent que les systèmes d'apprentissage accumulent une dette propre aux dépendances de
données, aux boucles de rétroaction, à la configuration et aux interactions entre
composants \cite{sculley2015debt}. Amershi et al. soulignent notamment la difficulté de
versionner les données, de réutiliser les modèles et de maintenir des composants dont le
comportement est entremêlé avec celui du système qui les consomme
\cite{amershi2019software}. Une dérive de distribution, une modification du parseur ou un
changement de schéma peut donc invalider un modèle sans modification de son code.

\subsection{Promotion, surveillance et retour arrière}

La préparation à la production exige des tests sur les données, les caractéristiques,
les modèles et l'infrastructure, puis une surveillance après déploiement. La grille de
Breck et al. regroupe précisément des contrôles de préparation et de monitoring pour
réduire la dette technique \cite{breck2017mltest}. Le NIST AI RMF ajoute une dimension de
gouvernance : rôles définis, documentation, mesure des risques, suivi continu et revue
périodique \cite{nist2023airmf}.

Une procédure minimale de promotion doit ainsi conserver :
\begin{enumerate}
    \item la version et la provenance des données, du parseur, des caractéristiques et du modèle ;
    \item un jeu de comparaison gelé et une règle de décision définie avant son évaluation ;
    \item les métriques prédictives, les coûts et les contrôles par sous-groupes ;
    \item une validation explicite, l'identité du décideur et l'artefact promu ;
    \item les conditions de surveillance, d'alerte, de retrait et de retour arrière.
\end{enumerate}

Un mode \textit{dry-run}, un registre ou une fonction de comparaison constituent des
briques de cette chaîne ; ils ne démontrent pas à eux seuls une gouvernance opérationnelle.
Dans Ariel Logminer, la mémoire persistante et les chemins de promotion fournissent une
infrastructure auditée par des tests. En l'absence de cycles répétés comprenant candidat,
décision, déploiement, surveillance et retour arrière, le positionnement doit rester celui
d'un \textbf{support de gouvernance expérimental}, non d'une plateforme MLOps validée ni
d'un apprentissage continu autonome.

